% Compile with pdfLaTeX.
\documentclass[12pt,a4paper]{article}
\usepackage[T1]{fontenc}
\usepackage[utf8]{inputenc}
\usepackage{mathptmx}
\usepackage[margin=1in]{geometry}
\pagestyle{empty}
\setlength{\parindent}{0pt}
\setlength{\parskip}{8pt}
\frenchspacing

\begin{document}
\raggedright
\textbf{Jingxian Liu}

Fields/Subfields: Labor Economics, Demographic Economics, Applied Econometrics

Keywords: Wage inequality, Intergenerational mobility, Demographic structure,
Population~aging, Fertility, Gender~inequality, Labor~supply, Household~consumption

Papers:

Autor, D.~H., Katz, L.~F., \& Kearney, M.~S. (2008).
  Trends in U.S. wage inequality: Revising the revisionists.
  \textit{The Review of Economics and Statistics, 90}(2), 300--323.

Card, D., Heining, J., \& Kline, P. (2013).
  Workplace heterogeneity and the rise of West German wage inequality.
  \textit{The Quarterly Journal of Economics, 128}(3), 967--1015.

Kleven, H., Landais, C., \& Søgaard, J.~E. (2019).
  Children and gender inequality: Evidence from Denmark.
  \textit{American Economic Journal: Applied Economics, 11}(4), 181--209.

Chetty, R., Hendren, N., Kline, P., \& Saez, E. (2014).
  Where is the land of opportunity? The geography of intergenerational
  mobility in the United States.
  \textit{The Quarterly Journal of Economics, 129}(4), 1553--1623.

Jaimovich, N., \& Siu, H.~E. (2009).
  The young, the old, and the restless: Demographics and business cycle volatility.
  \textit{American Economic Review, 99}(3), 804--826.
\end{document}
